\chapter{Opening}

This document is meant to do one job well:
create alignment on where Connector enters the market, what we sell first, and how we turn early customer interest into paid pilots.

It is intentionally simple.
It is not a pitch deck.
It is not a technical architecture paper.
It is not a full company plan.

It is a cofounder clarity document.

\begin{conceptbox}{Core statement}
We are not taking the full system to market.
We are entering through a narrow wedge where AI decisions need to be traceable and provable.

\smallskip
\centering
\Large\bfseries Connector = execution control layer for AI decisions.
\end{conceptbox}

\section{What this paper is trying to answer}

After reading this, a cofounder should be able to answer three simple questions:

\begin{itemize}
  \item Who do we call first?
  \item What do we sell first?
  \item Why will they pay?
\end{itemize}

If those answers are clear, we have something usable.
If those answers are still fuzzy, the strategy is not ready.

\section{What this paper is not trying to do}

This paper is not trying to explain every part of Connector.
It is not trying to carry the full technical depth of the platform.
And it is not trying to impress anyone with complexity.

We already have deeper material for that.
What we need here is a document that makes action easier.

\begin{warnbox}{Discipline point}
If the reader finishes this paper understanding the architecture but still cannot explain the entry wedge, then the paper failed.
\end{warnbox}

\section{Simple framing}

The market does not buy big visions first.
It buys specific relief for a specific pain.

For Connector, that pain is straightforward:
when an AI workflow takes action and something goes wrong, most teams cannot explain what happened, why it happened, or what evidence they can show.

That is where we enter.
